% Transcription — fonds Alexandre Grothendieck, shelfmark 161-5, batch 2 (pages 21-28).
% Established from the facsimile archives/batches/161-5/batch-02.pdf (a mirror of
% https://grothendieck.umontpellier.fr/161-5.pdf), per the skill
% transcribe-grothendieck. The batch file carries no cover sheet: PDF page N is
% archivists' page N+20, as the pencilled numbers bottom-left ("21" ... "28")
% confirm.
%
% Pass: Opus 5.5 (claude-opus-5-5), 2026-09-22 — first pass, unchecked against the pages by a human.
% Revised 2026-09-23 (Opus 5.5), after re-reading pp. 22-28 on the facsimile at
% 300-900 dpi: (1) the K/L exchange from p. 23 c) on is recorded as what the
% pages write — the p. 23 « sic » note, which treated it as an isolated slip,
% is replaced by one note flagging the exchange through p. 28, and the
% header says so below; letters in the body unchanged; (2) a note on p. 22 a)
% that the L of « L-vectoriel » is written in the interline; (3) « in two
% hands of the same work » replaced by « in three media ».
%
% What the batch is.
%   21     The last printed page (p. 50) of the offprint that ends batch 1: a
%          Collège de France course summary on p-divisible groups and Hodge-
%          Tate modules, ending with its « Publications » and « Missions »
%          lists. Printed, not his, and nothing on it is in his hand (the
%          marks in the right margin are show-through from the verso): absent,
%          following the offprint precedent of folder 22 batch 1.
%   22-28  Seven leaves of manuscript, his, in three media:
%          22-23 in ink on squared paper, a fair statement of the set-up;
%          24-26 in pencil on ruled paper, three successive attempts at the
%          list of data 1)-6) (p. 24 abandoned half-way, p. 25 restarting at
%          1), p. 26 continuing from 4)); 27-28 in pencil on squared paper,
%          the non-abelian cohomology of the extension (E) and a Question.
%          The lower half of p. 28 is a separate computation in ink (tori
%          over a p-adic field, a Néron model), unrelated to the pencil
%          argument above it; it is transcribed as a section of its own.
%
% Mathematics. K a local field of mixed characteristic with perfect residue
% field k, V its ring of integers, W ⊂ V the Witt ring, L = Frac W, Ω the
% completion of an algebraic closure, π = Gal(K̄/K) — so on p. 22 and
% p. 23 a)-b). From p. 23 c) to the end of p. 28 the pages use K and L the
% other way round, and consistently: T_DR(M) = T_cr(M) ⊗_K L (p. 23), Q_L
% « descendant Q_Ω » (p. 24), π = Aut_L(Ω) and Q_Ω « se descend à L »
% (p. 25), the « L-foncteur fibre filtré » DR_L, the descent « de L : K » and
% P_K ≃ P_K^(#) (p. 26), G'_S(S)^π = G'_L(L) (pp. 27-28) — L is there the
% base field, K the fraction field of W. The body keeps his letters as
% written; one \note at p. 23 c) flags the exchange, and the summary below
% also uses his letters page by page. The category M_K of
% Hodge-Tate Q_p-sheaves over K, its effective and « weight 0,1 » parts, the
% full embedding (Tate) of Barsotti-Tate groups up to isogeny over V into
% M_K^(1), and the Tannakian subcategory M^BT they generate. The fibre
% functor T_p, its group G (algebraic envelope of π), and the Hodge
% structure (Ω, π) → j^!: G_m → G^!_Ω, the torsors 𝔓, Q_Ω = 𝒯 ∧ 𝔓, their
% descent to K. For M ∈ BT: the crystalline L-space T_cr(M), its
% F-isocrystal structure, the filtration of T_DR(M) and the isomorphism
% gr T_DR ≃ T_Hdg — and the Question (p. 23): do these data extend to M^BT
% as ⊗-functors, uniquely? Pp. 24-26 axiomatise the answer as a list of
% data 1)-6) (group G over Q_p, ρ, j^!, the descent Q_L of Q_Ω, the torsor
% R_L under U''_L giving a filtered fibre functor DR_L, the descent datum
% L : K and the Frobenius-type isomorphism P_K ≃ P_K^(#)), recast in the end
% as an action of an extension Γ of Z containing π on P_Ω. Pp. 27-28: the
% sections of the resulting extension 1 → G'_S(S) → E → Γ → 1 classify the
% extensions of the torsor structure; Question: are those extending a given
% section over π conjugate under G'_S(S)^π = G'_L(L), i.e. is H^1(Γ, ·) →
% H^1(π, ·) « injective »?
%
% Notation fixed for the batch. His script M is \mathcal{M}; the circled P
% (a torsor under G_Ω with the diagonal π-action) is \mathfrak{P}, distinct
% from the plain P_L, P_K, P_Ω of p. 26; the script T (the G_m-torsor
% attached to T_p(Ω^*)) is \mathcal{T}; the contracted product written ∧
% with the group underneath is \wedge^{H}. G_m is \mathbb{G}_m. The symbol
% he writes as a curly S on pp. 27-28 (G'_S(S), (S, Γ; G_S)) is set as S;
% it is not defined in this batch. The small « g » of p. 28 (g_σ, g_K) is
% kept as g, and his « σ » there, an index for the ring of integers, as σ.
%
% Readings left open: the struck opening of 5) on p. 24 and the scribbled-
% out line of p. 25; the degrees of M_K^(n) on p. 22 (index read (n), the
% text says « 0 et 1 »); the superscript of P_K^(#) on p. 26; the sign
% « m » that labels rows on p. 28.
\documentclass[11pt,a4paper]{article}
% Le préambule commun à toutes les transcriptions du fonds Grothendieck.
%
% Il tient en une idée : **l'incertitude se compose, elle ne se résout pas.**
% Une transcription qui lisse un mot douteux a détruit la seule chose pour
% laquelle on la faisait — un lecteur doit pouvoir savoir, à chaque endroit,
% ce qui était sur le feuillet et ce qui a été ajouté. Les six macros qui
% suivent sont donc l'appareil critique, et non de la décoration.
%
% Les mêmes commandes sont reconnues par `scripts/render.mjs`, qui produit la
% vue de lecture du site. Une macro ajoutée ici doit l'être aussi là-bas,
% faute de quoi le rendu échouera bruyamment — ce qui est voulu : un rendu
% silencieusement faux est pire qu'un rendu absent.

\ProvidesPackage{grothendieck}[2026/08/08 Transcriptions du fonds Grothendieck]

\usepackage{amsmath,amssymb,amsthm}
\usepackage{mathrsfs}
\usepackage{stmaryrd}
% Les diagrammes commutatifs sont partout dans ce fonds : c'est la langue
% même de la géométrie algébrique de Grothendieck, et une transcription qui
% les rendrait en prose ne transcrirait rien.
\usepackage{tikz-cd}
% Français par défaut — c'est la langue des feuillets — mais l'anglais est
% chargé aussi : la traduction et le résumé sont des documents anglais, et la
% typographie française (espaces avant les deux-points) y serait une faute.
% Ces documents basculent par \selectlanguage{english} en tête de corps.
\usepackage[english,french]{babel}
\usepackage{fontspec}
\usepackage{geometry}
\geometry{a4paper,margin=25mm,marginparwidth=28mm}
\usepackage{marginnote}
\usepackage{xcolor}
% ulem, not soul. `soul`'s \st and \ul are built on a character-by-character
% scan that XeLaTeX's font machinery breaks: the document fails with an
% undefined control sequence rather than degrading. ulem does the same two jobs
% and is Unicode-engine safe. `normalem` stops it hijacking \emph, which the
% transcriptions use for Grothendieck's own emphasis.
\usepackage[normalem]{ulem}
\usepackage{hyperref}
\hypersetup{colorlinks=true,linkcolor=black,urlcolor=black,pdfborder={0 0 0}}

% --- Métadonnées du lot -----------------------------------------------------
% Reprises telles quelles par le script de rendu, qui les affiche en tête de
% la vue de lecture. Elles fixent ce que le fichier prétend couvrir : sans
% elles, une transcription flotte sans point d'attache dans un fonds de seize
% mille feuillets.

\newcommand{\folder}[1]{\gdef\grothFolder{#1}}
\newcommand{\batch}[1]{\gdef\grothBatch{#1}}
\newcommand{\leaves}[2]{\gdef\grothFirst{#1}\gdef\grothLast{#2}}
\newcommand{\foldertitle}[1]{\gdef\grothFoldertitle{#1}}
\gdef\grothFolder{?}\gdef\grothBatch{?}\gdef\grothFirst{?}\gdef\grothLast{?}\gdef\grothFoldertitle{}

% --- Le résumé --------------------------------------------------------------
%
% Il ouvre la lecture modernisée, sous le titre « Résumé », et il est composé
% en petit. Ce n'est pas un ornement : c'est le seul passage du document qui
% ne soit pas des mathématiques, et un lecteur doit pouvoir le reconnaître
% avant de l'avoir lu — puis le sauter s'il connaît déjà le sujet, ou s'y
% arrêter s'il ne le connaît pas. Le filet qui le suit marque la reprise du
% corps.

\newenvironment{resume}
  {\small\setlength{\parskip}{0.45em}}
  {\par\medskip\hrule\medskip}

% --- Mots-clés ---------------------------------------------------------------
%
% En anglais, en clôture du résumé de la lecture modernisée : le vocabulaire
% moderne sous lequel chercher ce que les feuillets construisent. Le manifeste
% du site (`npm run manifest`) les extrait pour étiqueter la cote entière —
% cette ligne est la source unique des tags, il n'y a pas de fichier à côté.
\newcommand{\keywords}[1]{\par\smallskip\noindent
  {\footnotesize\textsc{Keywords} --- \textit{#1}}\par}

% --- L'appareil critique ----------------------------------------------------

% Le numéro de feuillet, en marge. C'est l'ancre qui permet de régler un
% désaccord sur une formule en regardant *un* feuillet plutôt qu'une chemise
% de six cents. Le site s'en sert aussi pour tourner les pages du fac-similé
% quand on fait défiler la transcription.
\newcommand{\leaf}[1]{%
  \marginnote{\footnotesize\color{gray!70}\textsf{#1}}%
  \label{leaf:#1}\ignorespaces}

% Illisible : on ne devine pas. Un mot inventé qui se lit comme les autres est
% le pire résultat possible.
\newcommand{\ill}{\textcolor{red!60!black}{[\,\ldots\,]}}

% Lecture douteuse : le mot est donné, et signalé comme incertain.
\newcommand{\uncertain}[1]{\uline{#1}}

% Ajout de l'éditeur — une lettre manquante, un mot sous-entendu.
\newcommand{\add}[1]{\textcolor{blue!50!black}{[#1]}}

% Biffé par Grothendieck lui-même. Ce qu'il a rayé fait partie du document :
% une rature dit souvent où la pensée a changé de direction.
\newcommand{\struck}[1]{\sout{#1}}

% Note du transcripteur, sur l'état matériel du feuillet.
\newcommand{\note}[1]{\par\smallskip{\small\color{gray!60!black}\textit{[#1]}}\par\smallskip}

% Note marginale de Grothendieck, distincte du corps du texte.
\newcommand{\marginal}[1]{\par\smallskip\noindent\hspace{1em}%
  {\small\color{gray!60!black}#1}\par\smallskip}

% --- Titre ------------------------------------------------------------------

\AtBeginDocument{%
  \begin{center}
    {\large\bfseries Fonds Alexandre Grothendieck --- cote n\textsuperscript{o}~\grothFolder}\\[2pt]
    {\itshape\grothFoldertitle}\\[4pt]
    {\small feuillets \grothFirst--\grothLast\ (lot \grothBatch)}\\[2pt]
    {\footnotesize\color{gray!60!black}Transcription établie d'après le fac-similé numérisé
     par l'Université de Montpellier.\\
     Les lectures incertaines sont soulignées, les illisibles notées [\,\ldots\,],
     les ajouts entre crochets.}
  \end{center}
  \vspace{1em}}

\endinput


\folder{161-5}
\batch{2}
\pages{21}{28}
\dating{1967-1969 — le groupe « Dossiers rassemblés par Grothendieck sur différentes thématiques » (161-1 à 162-6) est daté [après 1961-vers 1977]}
\watermark{Édition de démonstration}
\foldertitle{Barsotti-Tate : tirés à part annotés (1967-1969), notes manuscrites (s.d.).}

\begin{document}

\subsection*{Notes manuscrites}

\page{22}
\struck{\ill{}} $K$ corps local d'inégales caractéristiques, corps résiduel $k$
de car.\ $p > 0$ ($k$ parfait!), anneau d'entiers $V$, sous-$p$-anneau $W$ de
corps résiduel $k$, corps des fractions $L$. On choisit une clôture algébrique
$\bar{K}$ de $K$, son complétion $\Omega$.

\[
\begin{array}{ccccccc}
L & \subset & K & \subset & \bar{K} & \subset & \Omega \\
\cup & & \cup & & \cup & & \cup \\
W & \subset & V & \subset & \bar{V} & \subset & \hat{\bar{V}}
\end{array}
\]
\note{ce tableau d'inclusions est écrit dans la marge de gauche, en regard des
cinq premières lignes.}

On pose
\[
\pi = \mathrm{Gal}(\bar{K}/K), \quad \pi' = \mathrm{Gal}(\bar{K}/L) \;\ldots
\]

$\mathcal{M}_K$ = catégorie des $\mathbb{Q}_p$-faisceaux du type Hodge-Tate sur
$K$ ($\otimes$-catégorie sur $\mathbb{Q}_p$)

$\mathcal{M}^{+}_K = \varinjlim \mathcal{M}^{(n)}$ ; ss-catégorie des objets
« effectifs » i.e.\ avec \struck{\uncertain{graduation}} $\mathrm{gr}(M_\Omega)$
à degrés positifs

\uncertain{$\mathcal{M}^{(n)}_K$} $= \{$ \struck{\ill{}} catégorie pleine de
$\mathcal{M}$ formée des $M$ avec $\mathrm{gr}\, M_\Omega$ à degrés $0$ et $1$
\note{l'exposant se lit « (n) », mais la condition « à degrés 0 et 1 » et la
ligne suivante, où figure $\mathcal{M}^{(1)}_K$, donnent plutôt
$\mathcal{M}^{(1)}_K$ ; l'accolade ouverte n'est pas refermée.}

$\widetilde{\mathrm{BT}}_V$ catégorie des BT à isog.\ près sur $V$
\note{l'indice $V$ est écrit sur un $K$ biffé.}

Foncteur pl.\ fid.\ (par Tate)
\[
\widetilde{\mathrm{BT}}_V \hookrightarrow \mathcal{M}^{(1)}_K
\]
\marginal{on identifie $\widetilde{\mathrm{BT}}_V$ à ss-catégorie pleine de
$\mathcal{M}^{(1)}_K$}

$\mathcal{M}^{\mathrm{BT}}_K \subset \mathcal{M}_K$ $\otimes$-\struck{\ill{}}
catégorie tannakienne engendrée par $\mathcal{M}^{(1)}_K$

$T^{\Omega}_p = T_p : \mathcal{M}_K \to$ \struck{\ill{}} $\mathrm{Modf}(\mathbb{Q}_p)$
\quad foncteur fibre sur $\Omega$

$G^{\Omega} = G = \mathrm{Aut}(T_p)$ \quad pro-groupe alg.\ sur $\mathbb{Q}_p$
\note{le second $G$ portait un exposant, biffé.}

\[
\pi \xrightarrow{\ \rho\ } G(\mathbb{Q}_p)
\]
$G$ est l'enveloppe algébrique de $\pi$ (NB $G/G^{\circ}$ non \ill{})

Pour $M \in \mathrm{Ob}\,\mathcal{M}$,
\[
G_M = \mathrm{Im}\bigl(G \to \underline{\mathrm{Aut}}_{\mathbb{Q}_p}(T_p(M))\bigr)
\subset \underline{\mathrm{Aut}}_{\mathbb{Q}_p}(T_p(M))
\]
\[
\bigl(\pi \xrightarrow{\ \rho_M\ } G_M(\mathbb{Q}_p) \quad \cdots\cdots \bigr)
\]

Structure sur $(\Omega, \pi)$ :
\quad $j^{!} : \mathbb{G}_{m,\Omega} \to G^{!}_\Omega$, \quad $\mathfrak{P}$
(avec section $e$ sur $\Omega$, \uncertain{non} invariante sous $\pi$), \quad
$G_\Omega$

\uncertain{id.} :
\quad $j_\Omega : \mathbb{G}_{m,\Omega} \to G^{\mathrm{Hdg}}_\Omega$, \quad
$Q_\Omega = \mathcal{T} \wedge^{\mathbb{G}_{m,\Omega}} \mathfrak{P}$, \quad
$G_\Omega$

\marginal{[$\mathcal{T}$ \uncertain{$= \mathcal{T}^{\Omega}$} \ill{}
$\mathbb{G}_{m,\Omega}$-tors.\ sur $(\Omega, \pi)$ associé à
$T_p(\Omega^{*}) \otimes_{\mathbb{Z}_p} \Omega$]}

par descente sur $K$ :
\quad $j_K : \mathbb{G}_{m,K} \to G^{\mathrm{Hdg}}_K$, \quad $Q_K$, \quad $G_K$,
\quad $T_{\mathrm{Hdg}}(M) = Q_K \wedge^{G} M_K$
\note{les trois lignes sont disposées en tableau : à gauche l'étiquette, puis
chaque fois un morphisme de $\mathbb{G}_m$ dessiné en diagonale montante, le
torseur, le groupe. L'exposant de $T_{\mathrm{Hdg}}$ est biffé.}

Pour $M \in \mathrm{Ob}\,\widetilde{\mathrm{BT}}_V$, on a
\struck{$T_{\mathrm{Hdg}}(M)$} \ill{}

a) Un $L$-vectoriel \struck{\uncertain{filtré}} $T_{\mathrm{cr}}(M)$
\struck{et un isom.\ $\mathrm{gr}\, T_{\mathrm{DR}}(M)_K \simeq T_{\mathrm{Hdg}}(M)$}
\note{le $L$ de « $L$-vectoriel » est écrit en interligne, d'un trait très
appuyé ; on ne voit pas s'il recouvre une autre lettre.}

\page{23}
b) Une \struck{donnée} structure de F-isocristal sur $T_{\mathrm{cr}}(M)$

c) Une filtration de $T_{\mathrm{DR}}(M) = T_{\mathrm{cr}}(M) \otimes_K L$
\note{à partir d'ici et jusqu'à la fin de la page 28, les lettres $K$ et $L$
sont employées à l'inverse de la page 22, où $L$ est le corps des fractions de
$W$, contenu dans $K$ : $T_{\mathrm{cr}}(M)$, dit en a) un $L$-vectoriel, est
ici tensorisé par $\otimes_K L$ ; page 24, $Q_L$ « descendant $Q_\Omega$ » ;
page 25, $\pi = \mathrm{Aut}_L(\Omega)$ et $Q_\Omega$ « se descend à $L$ » ;
page 26, le « $L$-foncteur fibre filtré », la descente « de $L : K$ » et
$P_K \simeq P^{(\sharp)}_K$ ; pages 27 et 28, $G'_S(S)^{\pi} = G'_L(L)$.
L'échange est suivi : $L$ y est le corps de base, $K$ le corps des fractions
de $W$. Les lettres sont transcrites telles qu'il les écrit.}

d) Un isomorphisme
$\mathrm{gr}(T_{\mathrm{DR}}(M)) \simeq T_{\mathrm{Hdg}}(M)$

ces données étant fonctorielles en $M \in \mathrm{Ob}\,\widetilde{\mathrm{BT}}_V$.

\emph{Question.} Peut-on prolonger ces données à $\mathcal{M}^{\mathrm{BT}}$,
\struck{\ill{} \uncertain{A-t-on}} avec $T_{\mathrm{cr}}$ un $\otimes$-foncteur,
la structure b) compatible avec $\otimes$ (et Hom), c) et d) aussi. A-t-on
unicité du $T_{\mathrm{cr}}$ prolongé, à isom.\ unique près ?

\page{24}
1) gr.\ alg.\ $\boxed{G}$ sur $\mathbb{Q}_p$

2) hom continu $\boxed{\rho : \pi \to G(\mathbb{Q}_p)}$

d'où $G^{!}_\Omega \;\text{—}\; \mathfrak{P}^{!}_\Omega \;\text{—}\; G_\Omega$

3) $\boxed{\mathbb{G}_{m,\Omega} \xrightarrow{\ j^{!}\ } G^{!}_\Omega}$
\quad d'où
\[
Q_\Omega = \mathcal{T} \wedge^{\mathbb{G}_{m,\Omega}} \mathfrak{P}^{!}_\Omega
\qquad G_\Omega
\]
$G''_\Omega = \mathcal{T}^{*} \wedge^{\mathbb{G}_{m,\Omega}} (G^{!}_\Omega,$
\uncertain{$\mathrm{int}$}$)$, \quad
$j_\Omega : \mathbb{G}_m \to G''_\Omega$
\note{le $\mathcal{T}$ est écrit par-dessus une autre lettre.}

\struck{4)}

4) $\boxed{Q_L \text{ sous } G_L}$ descendant $Q_\Omega$ sous $G_\Omega$
\marginal{unique}
\note{« unique » en marge de gauche, au bout d'une accolade qui embrasse 3) et
4).}

\[
U''_L \subset G''_L \;\text{—}\; Q_L \;\text{—}\; G_L, \qquad
j_L : \mathbb{G}_{m,L} \to G''_L
\]
\note{dans « $Q_L$ — $G_L$ », les indices $L$ sont écrits sur un $\Omega$.}

5) \struck{\uncertain{Un ensemble} $\boxed{\Delta}$ d'opérations $\rho$ de
$\Gamma$ ($\supset \pi$) sur $Q_\Omega$ \ill{} d'homomorphismes
$\Gamma \to G''_L(\Omega)$ \ill{}} prolongeant celle de $\pi$, comp.\ avec
celle de $\Gamma$ sur $\Omega$, \struck{tel que si $\rho$ est l'une, les autres
\ill{}} formant orbite sous $U''_L(L)$ opérant par (pour $u \in U''_L(L)$)
\[
(u.\rho)(\gamma) = \tau(u)\,\rho(\gamma)\,\tau(u)^{-1}
\]
\struck{$= (\mathrm{int}\,\tau(u)).\rho(\gamma)$}, i.e.\
$u.\rho = \mathrm{int}(\tau(u))\,\rho$.
\note{tout le passage qui va de « tel que » jusqu'ici est en outre annulé par
trois longs traits obliques ; la première ligne de 5) est biffée d'un trait
horizontal, et le début de la deuxième ne se lit pas.}

\struck{[NB si $u.\rho = \rho$ i.e.\ $\tau(u)$ et $\rho$ commutent :
$\tau(u)$ pour \ill{}}

Un torseur (\uncertain{commutatif}) $R_0$ à droite $\Delta$ sous $U''_L(L)$,
et un $U''_L(L)$-hom de
\[
\boxed{\rho : \Delta \to \text{Op.\ de } \Gamma \text{ sur } (\Omega, Q_L, G_L)}
\]
qui sont compatibles avec les op.\ qu'on a de $\Gamma$ sur $\Omega$, et celle
de $\pi \subset \Gamma$ sur $Q_L$… lorsque on fait opérer $U''_L(L)$ à
droite sur le deuxième ens.\ par $u.\alpha = \tau(\alpha)^{-1}\, u\, \tau(\alpha)$
\note{« à droite $\Delta$ » est ajouté en interligne au-dessus de $R_0$, qu'il
remplace peut-être : c'est $\Delta$ qui est la source de $\rho$.}

\page{25}
1) groupe alg.\ $\boxed{G}$ sur $\mathbb{Q}_p$. (Il dépend en fait du choix de
$\Omega$. cf plus bas 2°)

2) Hom $\boxed{\rho : \pi \to G(\mathbb{Q}_p)}$

[L'op.\ de $\pi = \mathrm{Aut}_L(\Omega)$ sur $G = G^{(\Omega)}$ est donnée par
$\gamma \mapsto \mathrm{int}(\rho(\gamma))$]

\struck{3)} On fait opérer $\pi$ sur $G_\Omega$ \emph{diagonalement}, soit
$G^{!}_\Omega$.

3) Hom $\boxed{j^{!} : \mathbb{G}_{m,\Omega} \to G^{!}_\Omega}$ \quad inv.\ par $\pi$

\marginal{\uncertain{$G_\Omega$}}

On appelle $S_\Omega(1)$ \struck{\ill{}} la \struck{\ill{}} $\Omega$-module
inversible $T_p(\Omega^{*}) \otimes_{\mathbb{Z}_p} \Omega$ avec op.\ diagonale
de $\pi$, \struck{\ill{}} et $\mathcal{T}_\Omega$ le
$\mathbb{G}_{m,\Omega}$-torseur associé. Posons
$Q_\Omega = \mathcal{T} \wedge \mathfrak{P}$ \struck{\ill{}} torseur (à dr.)\
sous $G_\Omega$, \ill{} commutant aux $G^{!}_\Omega$ tordu par
$(\mathcal{T}, j)$ soit $G''_\Omega$ [correspondant aux foncteurs fibres sur
$\mathrm{Rep}(G)$ (gradués)
\note{une ligne entière, entre « Posons » et « torseur », est raturée à la
mine jusqu'à l'illisible ; on y devine $\mathfrak{P}$, $Q$ et
$\mathbb{G}_{m,\Omega}$. La formule $Q_\Omega = \mathcal{T} \wedge \mathfrak{P}$
est écrite au-dessus, dans un cadre. « Posons » est en interligne.}

\[
\mathbb{G}_{m,\Omega} \xrightarrow{\ j_\Omega\ } G''_\Omega
\qquad
M \mapsto \sum \mathrm{gr}^{i}_{j}(M_\Omega) \otimes S_\Omega(-i)
= \mathrm{Hdg}(M)_\Omega
\]
\note{devant « Hdg » un trait vertical qui peut être un $T$ ou le crochet
fermant de la ligne précédente.}

\emph{Axiome.} Le \struck{torseur} $Q_\Omega$ \struck{\ill{}} se descend à $L$
en un $G_L$-torseur $Q_L$ (unique à isom.\ unique près) comme $G''_L$, et
$\mathbb{G}_{m,L} \xrightarrow{\ j_L\ } G''_L$.
\note{au-dessus du mot biffé « torseur », une surcharge entre parenthèses non
lue ; sous la ligne, $Q_L$ et $G_L$ sont repris dans un ovale.}

Correspond à fct fibre gradué $\mathrm{Hdg}_{j_L}$

\marginal{\emph{NB} Un $j$ satisfaisant aux conditions est unique. Signifie
que les $M_\Omega$ sont de Hodge. \ill{}}

\uncertain{2°} Il y a deux foncteurs fibres de $\mathrm{Rep}(G)$ dans
$\mathrm{Mod}(\Omega, \pi)$

a) $M \mapsto M_\Omega$ \quad opération \uncertain{par transport} des op.\
naturelles de $\pi$ sur $\Omega$

b) $M \mapsto M^{!}_\Omega$ \quad opération diagonale

--- foncteurs $T_{p,\Omega}$ et $T^{!}_{p,\Omega}$

\struck{\ill{}} On a donc un \struck{$(\Omega, \pi)$} torseur (à droite) sur
$(\Omega, \pi)$ du groupe $G_\Omega$, soit $\mathfrak{P}_\Omega$, donné par le
1-cocycle
\[
\rho : \pi \to G(\mathbb{Q}_p) \subset H^{0}(\Omega, G_\Omega) = G(\Omega).
\]
Le \struck{groupe} $\underline{\mathrm{Aut}}_{G_{\Omega,\pi}}(\mathfrak{P})$
\uncertain{$(\Omega, \pi)$} est $G^{!}_\Omega$.
\note{« 2° » est un chiffre cerclé en marge de gauche, lu d'après le renvoi
« cf plus bas 2° » du haut de la page.}

\page{26}
4) Soit $U''_L \subset G''_L$ \emph{défini} par $j_L$. On a un torseur (à
droite) \uncertain{canonique} (mais trivial) $\boxed{R_L \text{ sous } U''_L}$,
d'où $L$-foncteur fibre filtré
\[
\mathrm{DR}_L(M) = \bigl(R_L \wedge^{U''_L} Q_L\bigr) \wedge^{G_L} T_p(M)_L
\]
\note{une première écriture, $T_p(M)_L$, est biffée après le signe $=$.}

\struck{5)} Le foncteur fibre $\mathrm{DR}_L(M)$ correspond au $G_L$-tors.\
(à dr.)
\[
R_L \wedge^{U'_L} Q_L = P_L
\]

Soit $G'_L = \underline{\mathrm{Aut}}(P_L)$ \quad ($G'_L \simeq G_L$ non
canoniquement)

5) Structure de \emph{descente} de $L : K$ \emph{sur le torseur} $P_L$, d'où
\[
G''\qquad P_K \qquad G_K
\]

6) $G_K$-isom donné
\[
P_K \simeq P^{(\sharp)}_K
\]
\note{l'exposant de $P_K$ est un signe entre parenthèses qui tient du $\sharp$
et du $H$ ; il n'est pas lu avec certitude.}

On peut formuler ensemble 5) et 6) en regardant
\[
P_\Omega = R_L \wedge^{U''_\Omega} Q_\Omega = R_L \wedge^{U''_\Omega}
\]
\note{le premier membre de droite est surchargé (un $Q''$ biffé) ; le second
s'arrête après le signe $\wedge$.}

et on déclare que l'ext.\ $\Gamma$ de $\mathbb{Z}$ par
\uncertain{$\tilde{\pi} = \mathrm{Gal}(\bar{K}/\tilde{K})$} (qui contient
$\pi \subset \Gamma$) \emph{opère} sur $P_\Omega$ en \struck{\ill{}} prolongeant
l'op.\ qu'on a déjà sur $\pi$.

\page{27}
On a donc une \emph{extension}
\[
1 \to G'_S(S) \to \mathcal{E} \to \Gamma \to 1 . \qquad (\mathcal{E})
\]

\struck{Supposons que $\pi \subset \Gamma$ opère sur}

Les \emph{sections} de celle-ci sont les structures de torseur sur
$(S, \Gamma ; G_S)$ sur $G_S$ qui prolongent la structure de tor.\ sur
$(S, G_S)$. Si on a une telle section, on en déduit une opération de $\Gamma$
sur $G'_S(S)$ (qui en fait provient de son opération sur $G'_S$ lui-même !) et
les autres sections correspondent aux 1-cocycles de $\Gamma$ à valeurs dans
$G'_S(S)$. \struck{Les} L'ens.\ $\Gamma(\mathcal{E})$ des sections
\struck{forment} un ens.\ où $G'_S(S)$ opère (\uncertain{aucune restriction})
et deux \struck{\ill{}} $\Gamma(\mathcal{E})$ sont congrus ssi le 1-cocycle
faisant passer de l'un à l'autre est coh.\ à 1. Si on a un sous-groupe
$\pi \subset \Gamma$ et une section $\alpha_0 : \pi \to \mathcal{E}$ partielle,
on peut s'intéresser à l'\uncertain{ens.} des $g' \in G'_S(S)$ qui la fixent,
\ill{} des éléments de $G'_S(S)$ fixés par $\pi$ (y opérant \uncertain{comme on
pense}) $G'_S(S)^{\pi}$ $(= G'_L(L))$. Les solutions $\alpha$ qui prolongent
$\alpha_0$ sont donc telles que $G'_S(S)^{\pi}$ y opère. Si on a
\note{« est coh.\ à 1 » est ajouté en marge de gauche avec un trait de renvoi.}

\page{28}
un tel $\alpha$, les autres $\beta$ s'en déduisent par les 1-cocycles qui sont
\emph{triviaux sur} $\pi$.
\[
\begin{cases}
\varphi(\gamma\gamma') = \varphi(\gamma)\; {}^{\gamma}\varphi(\gamma') \\
\varphi(\pi) = \{e\}
\end{cases}
\]

\emph{Question.} Est-il vrai que les $\alpha$ qui prolongent $\alpha_0$ sont
conj.\ i.e.\ conjugués sous $G'_S(S)$ ? Il suffirait de savoir que (pour une
telle section fixée)
\[
H^{1}(\Gamma, G'_S(S)) \to H^{1}(\pi, G'_S(S))
\]
est injectif (ou plutôt « injectif » i.e.\ noyau unité). Si oui, les solutions
cherchées forment un \struck{\ill{}} espace homogène sous
$G'_S(S)^{\pi}$ $(= G'_L(L))$
\note{un long trait vertical à la mine traverse la page de haut en bas, et un
zigzag le longe dans la moitié inférieure ; ils séparent plutôt qu'ils ne
biffent, le texte qu'ils croisent étant repris tel quel. Trois traits
verticaux en marge de gauche marquent la Question.}

\subsection*{Tores sur un corps $p$-adique}
\note{la moitié inférieure de la page 28 est à l'encre, sous un trait
horizontal et la lettre H au crayon ; c'est un calcul indépendant de ce qui
précède. Les deux diagrammes sont côte à côte sur la feuille.}

$\mathbb{Q}_p$ --- $K$

$T_K$ \struck{$(K)$} $= \prod_{K/\mathbb{Q}_p} \mathbb{G}_m$,
\qquad $T_K(\mathbb{Q}_p) \simeq K^{*}$

\begin{tikzcd}[column sep=small, row sep=small, nodes={font=\scriptsize}]
 & \prod_{\sigma/\mathbb{Z}_p} \mathbb{G}_m \arrow[r] & T_{\sigma_K} \arrow[r] & \mathbb{Z} \arrow[r] & 0 \\
0 \arrow[r] & g_\sigma \arrow[r] & N \arrow[r] & \mathbb{Z} \arrow[r] & 0 \\
0 \arrow[r] & g_\sigma \arrow[u, Rightarrow] \arrow[r] & \tilde{N} \arrow[u] \arrow[r] & \mathbb{Z}_\sigma \arrow[u] \arrow[r] & 0 \\
0 \arrow[r] & g_K \arrow[r] & \tilde{N}_K \arrow[r] & \mathbb{Z}_K \arrow[r] & 0 \\
 & & g_K \times \mathbb{Z}_K \arrow[u, no head, "\simeq" description] & &
\end{tikzcd}

\marginal{modèle de Néron de $g_K$}
\note{la note marginale pointe par un trait sur $N$. La dernière ligne est
étiquetée à gauche « $\otimes_\sigma K$ », avec en dessous « $\| \; K$ » ; un
signe ressemblant à « m », non lu, précède la troisième et la quatrième ligne.}

\begin{tikzcd}[column sep=small, row sep=small, nodes={font=\scriptsize}]
 & 0 & \\
 & \mathbb{Z} \arrow[u] & \\
 & K^{*} \arrow[u] & \\
 & \sigma^{*}_K \arrow[u] & \\
 & 0 \arrow[u] & \\
\Gamma(N) \arrow[rr, "\simeq"] & & \Gamma(g) = K^{*} \arrow[ddll, "\pi"'] \\
\Gamma(\tilde{N}) \arrow[u, "2"'] \arrow[d, hook] & & \\
\Gamma(\tilde{N}_K) \simeq K^{*} \times \mathbb{Z} & &
\end{tikzcd}

\note{le $K^{*}$ de la colonne est celui de $T_K(\mathbb{Q}_p) \simeq K^{*}$
écrit plus haut ; l'isomorphisme $\Gamma(\tilde{N}_K) \simeq K^{*} \times
\mathbb{Z}$ est écrit verticalement sous $\Gamma(\tilde{N}_K)$, et la flèche
oblique, marquée $\pi$ — ici sans doute une uniformisante —, part de $K^{*}$.
Le « 2 » sur la flèche montante est une
lecture douteuse. À droite de la colonne, séparé d'elle par le même signe
« m », un second $\mathbb{Z}$ flanqué de flèches montantes, sans source
écrite.}

\end{document}
